\chapter{Exercises, counterexamples, and research projects}
The following problems are original exercises for the CON\textsuperscript{2}PHYS
curriculum. The objective is to prove each mathematical claim first
and only then write MATLAB experiments.

\begin{exercise}[Relations and functions]
Define $R\subseteq U\times T$ for two units firing at the same
timestamp. Prove that $R$ need not be a function $U\to T$.
Construct the corresponding function into $\mathcal P(T)$.
\end{exercise}
\begin{exercise}[Sampling failure]
Sample a 90-Hz sine wave at 100 Hz and identify a lower-frequency
sinusoid with the same samples. Explain why a digital filter applied
after sampling cannot recover the true generating frequency.
\end{exercise}
\begin{exercise}[Hilbert envelope]
Construct $x(t)=(1+0.3\cos 2\pi 2t)\cos 2\pi 40t$.
Calculate the analytic envelope and compare against a signal
whose envelope frequencies overlap its carrier spectrum.
Explain when the Bedrosian assumption fails.
\end{exercise}
\begin{exercise}[Spectral concentration]
Generate equal-length noisy signals with an off-bin sinusoid.
Compare ordinary periodograms, Welch, and DPSS multitaper
estimates using repeated simulations. Vary time-bandwidth
product and report bias versus variance.
\end{exercise}
\begin{exercise}[Principal subspace perturbations]
Generate a diagonal covariance with a controllable eigengap.
Perturb it by symmetric noise, compute leading subspace angles,
and compare their variation with the Davis--Kahan scaling.
\end{exercise}
\begin{exercise}[Random-matrix nulls]
Simulate isotropic Gaussian matrices with independent samples,
then introduce temporal autocorrelation and heterogeneous
channel noise. Compare each empirical eigenvalue histogram with
the corresponding idealized Marcenko--Pastur curve.
\end{exercise}
\begin{exercise}[Latent gain]
Generate counts with $N\mid G\sim\Poisson(\mu G)$, where
$G$ is a positive random variable of mean one.
Verify the predicted quadratic mean--variance relation and
test how a purely Poisson GLM misestimates uncertainty.
\end{exercise}
\begin{exercise}[Conditional point-process validation]
Fit a simulated point-process intensity with and without spike-history
effects. Examine rescaled inter-event intervals. Explain why
correctly transformed residuals do not identify a unique mechanism.
\end{exercise}
\begin{exercise}[Shared-input Granger counterexample]
Simulate an unobserved driver of two time series with different lags.
Show that adding measurements of the common source can change
inferred predictive direction or strength. Avoid interpreting
this as proof of synaptic connectivity.
\end{exercise}
\begin{exercise}[Inverse problem]
Construct a lead-field matrix with a nontrivial nullspace.
Find two distinct current patterns with identical electrode
predictions, then compare Tikhonov solutions for different
regularization strengths.
\end{exercise}

\section{Suggested cumulative projects}
Project I: from set membership to a reliable spike-count
representation. Project II: from convolution to LFP spectral
density and multitaper estimators. Project III: from covariance
to PCA stability and manifold caveats. Project IV: from
conditional probability to spike-train GLM validation.
Project V: from state-space observability to limits of
causal and mechanistic inference.
% --- expanded edition (October 2026) ---
\section{Further problems by chapter}
\begin{exercise}[Partitions and sorting]
Two spike-sorting runs on the same data give partitions $\mathcal P$ and
$\mathcal Q$. Prove that the number of spike pairs placed together by
both runs equals $\sum_{P,Q}\binom{|P\cap Q|}{2}$, and use this to define
the Rand index.
\end{exercise}
\begin{exercise}[Contraction and recurrent rates]
For $r=\tanh(wr+I)$, find the range of $w$ for which the fixed-point
iteration is guaranteed to converge for every $I$.
\end{exercise}
\begin{exercise}[PCA of a random walk]
Simulate 100 independent random walks of length 1000, compute the PCA
across time, and compare the leading components with the
Karhunen--Lo\`eve eigenfunctions of Brownian motion.
\end{exercise}
\begin{exercise}[Delta method for coherence]
Using the Fisher transformation $\operatorname{atanh}|\hat c|$, derive an
approximate confidence interval for coherence from $K$ tapers and
compare it with a jackknife over tapers in simulation.
\end{exercise}
\begin{exercise}[GHK inversion]
Given reversal potentials measured under two different extracellular
sodium concentrations, show that the ratio $P_{\rm Na}/P_{\rm K}$ is
identifiable from the GHK voltage equation, while the absolute
permeabilities are not.
\end{exercise}
\begin{exercise}[Causal phase estimation]
Fit an AR model to the past 500\,ms of an LFP and forecast 100\,ms ahead
to estimate the current theta phase. Compare the error with that of a
causal Butterworth filter and with the zero-phase reference.
\end{exercise}
\begin{exercise}[Hawkes overdispersion]
Simulate a Hawkes process with exponential kernel and branching ratio
$0.5$, and verify the long-window Fano factor $1/(1-\nu)^2$.
\end{exercise}
\begin{exercise}[Collider bias in event-triggered analysis]
Simulate two independent Poisson neurons, define population events as
bins where the summed count exceeds a threshold, and compute their
correlation within and outside events.
\end{exercise}
\begin{exercise}[Spiked covariance]
Plant a rank-one component of variance $1+\theta$ in $N=100$ isotropic
Gaussian neurons. For $T\in\{50,100,400,1600\}$, estimate the leading
sample eigenvalue and the squared alignment, and compare with the BBP
predictions.
\end{exercise}
\begin{exercise}[Entrainment tongue]
Drive a noisy phase oscillator with natural frequency 8\,Hz at
frequencies from 6 to 10\,Hz and several intensities. Map the region in
which the phase-locking value to the drive exceeds 0.8, and compare with
the Adler prediction.
\end{exercise}
\begin{exercise}[Entropy bias]
For a uniform distribution on $K=16$ symbols, estimate the plug-in
entropy bias for $T=20,50,200$ by simulation and compare with
$(K-1)/(2T)$.
\end{exercise}
\begin{exercise}[Time rescaling with a wrong model]
Simulate an inhomogeneous Poisson process with sinusoidal rate, fit a
constant-rate model, and show how the KS plot of rescaled intervals
reveals the misspecification.
\end{exercise}
